\chapter{Introducción para los jugadores}

\begin{CajaSecreto}
\lead{Nota para el Director:} lee el siguiente texto a tus jugadores antes de comenzar la primera sesión. Sitúa el tono en un otoño húmedo de 1992 en la margen izquierda de la ría de Bilbao: bicicletas, humo industrial, canciones de casete y rumores de desapariciones en el telediario.
\end{CajaSecreto}

\section{Lectura para la mesa}

\kobfigure{01_parque_markonzaga.jpg}{Tarde de otoño en el parque de Markonzaga: entre el juego y el humo de las fábricas.}

\begin{quote}
«Bienvenidos a Sestao, otoño de 1992. Mientras en Barcelona se apagan los ecos de las Olimpiadas y en Sevilla recogen los pabellones de la Expo, aquí en la margen izquierda de la ría del Nervión la realidad tiene otro color: el color del óxido, del acero mojado y del carbón.

La otrora boyante industria siderúrgica y naval que dio de comer a miles de familias se desmorona día a día víctima de la reconversión industrial. Frente a la ría se alzan los esqueletos de Altos Hornos de Vizcaya (AHV), escupiendo sus últimos estertores de humo sobre el agua marrón. Hacia el interior, las inmensas naves de Babcock \& Wilcox (B\&W) recuerdan que en los ochenta fabricaron los componentes de las centrales nucleares, incluida la malograda central de Lemóniz. En las esquinas y plazas aún se ven las huellas de las huelgas obreras, los parados que se reúnen en las tabernas y las secuelas de los duros años ochenta.

Pero vosotros tenéis entre once y quince años. Vuestro mundo son las bicicletas BH y Orbea California, las partidas de cinco duros en los futbolines y recreativos de Txetxu en la Gran Vía, el olor a regaliz negro y pipas Facundo, y las carreras bajando las cuestas hacia La Benedicta. El Sestao Sport lucha por no descender en Segunda, y los fines de semana subís en el funicular de La Reineta a los lagos de La Arboleda a buscar fósiles o explorar viejas galerías mineras.

Sin embargo, hay algo en el aire de este otoño que no huele solo a fundición. En las noches despejadas, sobre el agua de la ría brillan destellos esmeralda que ningún pescador sabe explicar. En los televisores de tubo, la imagen de Paco Lobatón en \emph{Quién sabe dónde} recuerda el rostro de Alex Etxebarria, vuestro amigo del colegio que desapareció sin dejar rastro en las cuevas de La Arboleda. Y hoy, al salir del colegio, una columna de humo negro y el aullido histérico de las sirenas de bomberos rasgan el cielo sobre las vallas electrificadas de Babcock \& Wilcox.

Ajustad las cadenas de las bicis y encended las linternas. Algo ha despertado bajo el asfalto de Sestao.»
\end{quote}

\section{Prólogo audiovisual: Entre barcos y acero}

Guion cinemático de un minuto y medio, pensado como prólogo inmersivo para proyectar o narrar antes de la Escena~1.

\begin{CajaCinematica}{Ficha técnica del prólogo}
\cine{Título:} SESTAO 1992: Entre barcos y acero\par
\cine{Duración:} 1 minuto y 20 segundos\par
\cine{Tono visual:} Nostalgia noventera, realismo obrero, lluvia incesante, luces de neón parpadeantes en los recreativos frente a las chimeneas titanio de AHV y el brillo verde anómalo del Éter.\par
\cine{Banda sonora:} Sonido analógico de sintetizador lo-fi, chirrido metálico de grúas de astillero, sirenas lejanas de bomberos y el golpeteo del sirimiri sobre cuadros de bicicleta.
\end{CajaCinematica}

\begin{CajaCinematica}{Escena 1 -- La Ría de óxido (0:00 -- 0:20)}
\cine{Visual:} Plano general aéreo de la ría del Nervión al atardecer bajo un cielo encapotado de nubes plomizas. Las chimeneas de Altos Hornos de Vizcaya recortadas contra el horizonte. A ras de suelo, un grupo de chavales en bicicletas pedalea con fuerza por el bidegorri de La Benedicta.\par
\cine{Efectos de sonido:} El traqueteo de piñones de bicicleta, gaviotas sobre el agua fangosa y el eco distante de sirenas industriales.\par
\cine{Voz en off (tono de adolescente nostálgico pero inquieto):}\par
\textit{«En 1992 decían que España era moderna. Hablaban de Barcelona, de Sevilla y del futuro. Pero en Sestao el futuro seguía oliendo a chatarra, a lluvia ácida y a humo de altos hornos...»}
\end{CajaCinematica}

\begin{CajaCinematica}{Escena 2 -- El pueblo y los recreativos (0:20 -- 0:40)}
\cine{Visual:} Ráfagas rápidas de la vida cotidiana: los bares de poteo, señores mayores con txapela leyendo el periódico con noticias de reconversión, y la entrada de la sala de recreativos de Txetxu iluminada por pantallas de Street Fighter y humo de tabaco. Un cartel de búsqueda desgastado pegado a una farola: \emph{«Desaparecido: Alex Etxebarria, 11 años»}.\par
\cine{Efectos de sonido:} Monedas de 25 pesetas cayendo, música chiptune distorsionada y el chasquido de un walkman encendiéndose.\par
\cine{Voz en off:}\par
\textit{«Pasábamos las tardes contando cinco duros en los futbolines y pedaleando hasta reventar. Intentábamos olvidar que hacía tres años Alex entró en las minas de La Arboleda... y nunca volvió a salir».}
\end{CajaCinematica}

\begin{CajaCinematica}{Escena 3 -- La explosión en B\&W (0:40 -- 1:00)}
\cine{Visual:} De repente, un relámpago verde esmeralda surge de las naves interiores de Babcock \& Wilcox. Una columna de humo negro se eleva mientras camiones de bomberos y coches de Protección Civil cruzan las barreras de seguridad con sirenas azulosas.\par
\cine{Efectos de sonido:} Una detonación sorda, zumbido eléctrico en alta frecuencia y el ulular estridente de sirenas.\par
\cine{Voz en off:}\par
\textit{«Pero hoy la sirena no avisaba de un turno de fundición. Algo reventó dentro de los laboratorios blindados de Babcock \& Wilcox. Algo que no era de este mundo...»}
\end{CajaCinematica}

\begin{CajaCinematica}{Escena 4 -- El despertar (1:00 -- 1:20)}
\cine{Visual:} En la penumbra de una huerta abandonada, una pequeña figura con impermeable empapado alza la mirada. Sus ojos reflejan un destello antinatural de luz verde cian. Las farolas de la calle parpadean y revientan a su paso.\par
\cine{Texto en pantalla:}\par
\textbf{ENTRE BARCOS Y ACERO: SESTAO 1992}\par
\textbf{Una aventura para Kids on Bikes}\par
\cine{Voz en off (susurro infantil entre estática de radio):}\par
\textit{«...el eclipse... el eclipse lo va a despertar todo...»}
\end{CajaCinematica}
